% Ledger sources: Q6604 (original near-wall proof), Q6630 (hostile audit,
% with the nu_D prose typo corrected), and Q6657 (independent audit in the
% current one-based convention).  The only external inputs used below are
% R_h<=3(h-1), Theorem~\ref{thm:sh-sigma-half}, and the previously proved
% interval multiplicity estimate m_I(v)<=1+4|I|^{2/3}.

\subsection{A near-wall estimate for close pairs of gaps}
\label{sec:nw-near-wall}

We retain the conventions of Sections~\ref{sec:orbit-energy}
and~\ref{sec:sigma-half}: $p\ge7$ is prime,
$I_p=\{1,\ldots,M\}$ with $M=p-2$, and
$\pi(n)=[b_n:c_n]\in\mathbb P^1(\mathbb F_p)$.
All counts below are nonwrapping.  Thus a condition involving
$\pi(r+h)$ always includes $r+h\le M$.  For convenience set $R_h=0$
when $h>M-1$.

For $1\le a<b$ define
\[
 J_p(a,b)
 =\#\{r\in I_p:r+b\le M,
          \ \pi(r)=\pi(r+a)=\pi(r+b)\}.
\]
For integers $1\le D\le H$ put
\begin{equation}\label{eq:nw-K-def}
 K_p(H,D)
 =\sum_{1\le d\le D}\ \sum_{H<h\le2H-d}J_p(h,h+d).
\end{equation}
Thus $K_p(H,D)$ counts same-base pairs of collision gaps whose smaller
gap lies in $(H,2H]$ and whose separation is at most $D$.

For $d\ge1$ and $u\in I_p$ define
\[
 Z_d=\{u\in I_p:u+d\le M,\ \pi(u)=\pi(u+d)\},
\]
\[
 \nu_D(u)
 =\#\{1\le d\le D:u\in Z_d\},
 \qquad
 S_1(D)=\sum_{1\le d\le D}R_d=\sum_{u\in I_p}\nu_D(u).
\]
The ordinary same-base pair count is
\begin{equation}\label{eq:nw-Q-def}
 Q_p(D)=\sum_{u\in I_p}\binom{\nu_D(u)}2.
\end{equation}
For the backward window define
\[
 W_{H,d}(u)
 =\#\{h:H<h\le2H-d,\ h\le u-1,
          \ \pi(u-h)=\pi(u)\},
\]
\begin{equation}\label{eq:nw-W-def}
 W_H(u)
 =\#\{h:H<h\le2H,\ h\le u-1,
          \ \pi(u-h)=\pi(u)\}.
\end{equation}
The endpoint condition $h\le u-1$ is the one-based nonwrapping condition
$u-h\ge1$.

\begin{lemma}[forward-window second moment]
\label{thm:nw-nu-identity}
For every integer $D\ge1$,
\begin{equation}\label{eq:nw-nu-identity}
 \sum_{u\in I_p}\nu_D(u)^2=S_1(D)+2Q_p(D).
\end{equation}
Consequently,
\begin{equation}\label{eq:nw-nu-L2}
 \sum_{u\in I_p}\nu_D(u)^2\le15D^{5/2}.
\end{equation}
\end{lemma}

\begin{proof}
The pointwise identity $x^2=x+2\binom x2$, summed over $u$, gives
\eqref{eq:nw-nu-identity}.  Here an unordered pair $d_1<d_2$ counted by
$\binom{\nu_D(u)}2$ is exactly the same-base collision pair
\[
 \pi(u)=\pi(u+d_1)=\pi(u+d_2)
\]
with gaps $d_1,d_2\le D$.  No shift of the base point is involved.

By~\eqref{eq:sh-onegap},
\[
 S_1(D)\le3\sum_{d=1}^{D}(d-1)
          =\frac32D(D-1)\le\frac32D^{5/2}.
\]
Theorem~\ref{thm:sh-sigma-half} gives
$Q_p(D)\le\frac{27}{4}D^{5/2}$, since its truncated gap parameter is at
most $D$.  Substitution in~\eqref{eq:nw-nu-identity} yields
\[
 \sum_u\nu_D(u)^2
 \le\left(\frac32+\frac{27}{2}\right)D^{5/2}
 =15D^{5/2}.
\]
\end{proof}

\begin{lemma}[backward-window injection]
\label{thm:nw-w-injection}
For every integer $H\ge1$,
\begin{equation}\label{eq:nw-w-injection}
 \sum_{u\in I_p}\binom{W_H(u)}2\le Q_p(2H).
\end{equation}
Moreover,
\begin{equation}\label{eq:nw-W-L2}
 \sum_{u\in I_p}W_H(u)^2
 \le\left(\frac92+54\sqrt2\right)H^{5/2}.
\end{equation}
\end{lemma}

\begin{proof}
An object counted by $\binom{W_H(u)}2$ is a uniquely ordered pair
\[
 H<h_1<h_2\le2H,
 \qquad h_2\le u-1,
 \qquad
 \pi(u-h_1)=\pi(u)=\pi(u-h_2).
\]
Set
\[
 r=u-h_2,
 \qquad
 g_1=h_2-h_1,
 \qquad
 g_2=h_2.
\]
Then
\[
 r\ge1,
 \qquad
 1\le g_1\le H-1,
 \qquad
 H<g_2\le2H,
 \qquad
 g_1<g_2,
 \qquad
 r+g_2=u\le M.
\]
Furthermore
\[
 r+g_1=u-h_1,
 \qquad
 r+g_2=u,
\]
so
\[
 \pi(r)=\pi(r+g_1)=\pi(r+g_2).
\]
Thus $(r,g_1,g_2)$ is a valid object counted by $Q_p(2H)$.  The map is
injective because its image recovers the source through
\[
 u=r+g_2,
 \qquad
 h_2=g_2,
 \qquad
 h_1=g_2-g_1.
\]
This proves~\eqref{eq:nw-w-injection}, including the ordered-pair and
nonwrapping conventions.

The first moment is
\[
 \sum_{u\in I_p}W_H(u)
 =\sum_{H<h\le2H}R_h
 \le3\sum_{h=H+1}^{2H}(h-1)
 =\frac{9H^2-3H}{2}
 \le\frac92H^2.
\]
Using $x^2=x+2\binom x2$, the injection, and
Theorem~\ref{thm:sh-sigma-half},
\begin{align*}
 \sum_uW_H(u)^2
 &\le \frac92H^2+2Q_p(2H)\\
 &\le \frac92H^2
      +2\cdot\frac{27}{4}(2H)^{5/2}\\
 &=\frac92H^2+54\sqrt2\,H^{5/2}\\
 &\le\left(\frac92+54\sqrt2\right)H^{5/2},
\end{align*}
which is~\eqref{eq:nw-W-L2}.
\end{proof}

\begin{theorem}[near-wall bound with $\eta=3/7$]
\label{thm:nw-near-wall}
For every prime $p\ge7$ and all integers $1\le D\le H$,
\begin{equation}\label{eq:nw-main-explicit}
 K_p(H,D)
 \le C_{\rm nw}\,H D^{11/7},
 \qquad
 C_{\rm nw}
 =\sqrt{15\left(\frac92+54\sqrt2\right)}.
\end{equation}
In particular,
\[
 K_p(H,D)\ll H D^{11/7}=H D^{2-3/7}
\]
uniformly in $p,H,D$, with no factor $p^\varepsilon$.
\end{theorem}

\begin{proof}
We first rewrite the count at its midpoint.  A term of
$K_p(H,D)$ consists of $(r,h,d)$ satisfying
\[
 1\le d\le D,
 \qquad
 H<h\le2H-d,
 \qquad
 r+h+d\le M,
\]
and
\[
 \pi(r)=\pi(r+h)=\pi(r+h+d).
\]
Put $u=r+h$.  Then
\[
 r=u-h,
 \qquad
 r+h+d=u+d,
\]
and the event becomes
\[
 \pi(u-h)=\pi(u)=\pi(u+d).
\]
The boundary conditions become $u+d\le M$ and $h\le u-1$.  Conversely,
these conditions recover the unique base $r=u-h$.  Hence
\begin{equation}\label{eq:nw-midpoint-identity}
 K_p(H,D)
 =\sum_{1\le d\le D}\ \sum_{u\in Z_d}W_{H,d}(u).
\end{equation}
Since all summands are nonnegative and
$W_{H,d}(u)\le W_H(u)$,
\begin{equation}\label{eq:nw-midpoint-envelope}
 K_p(H,D)\le\sum_{u\in I_p}\nu_D(u)W_H(u).
\end{equation}

Cauchy--Schwarz, Lemma~\ref{thm:nw-nu-identity}, and
Lemma~\ref{thm:nw-w-injection} give
\begin{align}
 K_p(H,D)
 &\le
 \left(\sum_u\nu_D(u)^2\right)^{1/2}
 \left(\sum_uW_H(u)^2\right)^{1/2}\notag\\
 &\le
 \sqrt{15\left(\frac92+54\sqrt2\right)}
 H^{5/4}D^{5/4}\notag\\
 &=C_{\rm nw}H^{5/4}D^{5/4}.
 \label{eq:nw-cauchy-branch}
\end{align}

There is also a pointwise branch.  For fixed $u$, all positions counted by
$W_H(u)$ lie in the interval
\[
 [\max(1,u-2H),\,u-H-1],
\]
which contains at most $H$ integers, all of projective color $\pi(u)$.
The previously proved interval multiplicity estimate therefore gives
\[
 W_H(u)\le1+4H^{2/3}\le5H^{2/3}.
\]
Together with~\eqref{eq:nw-midpoint-envelope} and
$S_1(D)\le\frac32D(D-1)$, this yields
\begin{equation}\label{eq:nw-pointwise-branch}
 K_p(H,D)
 \le\left(1+4H^{2/3}\right)S_1(D)
 \le\frac{15}{2}H^{2/3}D^2.
\end{equation}

The two powers cross at $D=H^{7/9}$.  If $D\le H^{7/9}$, then
$D^{3/7}\le H^{1/3}$, and~\eqref{eq:nw-pointwise-branch} gives
\[
 K_p(H,D)
 \le\frac{15}{2}H^{2/3}D^2
 \le\frac{15}{2}H D^{11/7}.
\]
If $D\ge H^{7/9}$, then $D^{9/28}\ge H^{1/4}$, and
\eqref{eq:nw-cauchy-branch} gives
\[
 K_p(H,D)
 \le C_{\rm nw}H^{5/4}D^{5/4}
 \le C_{\rm nw}H D^{11/7}.
\]
Since $C_{\rm nw}>15/2$, the two cases prove
\eqref{eq:nw-main-explicit}.
\end{proof}

\begin{remark}[position on the wall map]
\label{rem:nw-wall-map}
The proof retains the sharper two-branch estimate
\begin{equation}\label{eq:nw-sharp-min}
 K_p(H,D)
 \le\min\left\{
 \frac{15}{2}H^{2/3}D^2,
 \ C_{\rm nw}H^{5/4}D^{5/4}
 \right\}.
\end{equation}
The theorem packages this minimum into the near-wall form
$H D^{2-\eta}$ with the unconditional gain $\eta=3/7$.

There are also the global containments
\[
 K_p(H,D)\le Q_p(2H)
 \le\frac{27}{4}(2H)^{5/2}
 =27\sqrt2\,H^{5/2},
\]
and, from the abstract $H^2\log H$ theorem with $R_d\le3d$,
\[
 K_p(H,D)\le Q_p(2H)
 \le66(2H)^2\bigl(1+\log(2H)\bigr)
 =264H^2\bigl(1+\log(2H)\bigr).
\]
Thus the near-wall theorem is a local strip estimate, not a replacement
for the sharp global bound.  At the power scale $D=H^\alpha$, its exponent
is $1+11\alpha/7$.  It improves on the $H^{5/2}$ global exponent for
$\alpha<21/22$, and it gives a power saving over $H^2\log H$ for
$\alpha<7/11$; at $\alpha=7/11$ it gives $O(H^2)$ and removes the global
logarithm.  For larger $\alpha$, the $H^2\log H$ theorem is the stronger
global estimate.  The new information is therefore concentrated exactly
in the near-diagonal regime where two collision gaps are close.
\end{remark}
